\section{Problem Formulation}
\label{sec:problem}

\subsection{Transit, transfer and transition sets}
A robot with configuration $x\in\mathcal{X}\subset\R^{n_r}$ (arm and fingers)
manipulates an object with configuration $\alpha$ in a known static scene. The
state $q=(x,\alpha)$ lives in the composite configuration space
$\CS$~\cite{simeon2004manipulation}, and $\CSfree\subseteq\CS$ is the set of
states without collisions other than the contacts a grasp requires. The object
cannot move on its own. In the \emph{placement space} $\CP$ it rests stably,
and in the \emph{grasp space} $\CG$ the robot holds it. A manipulation path
alternates \emph{transit} segments in $\CP$, where the robot moves alone and
$\alpha$ is constant, and \emph{transfer} segments in $\CG$, where the robot
moves the object at a fixed relative pose. The two can only meet in the
\emph{transition sets} $\CG\cap\CP\cap\CSfree$, where the robot can grasp or
release the object. These are the nodes of the manipulation
graph~\cite{simeon2004manipulation}.

\subsection{Multi-regrasp tasks with a coupled object coordinate}
We study the tasks captured by the following assumption, which include caps,
knobs, valves and cranks.
\begin{assumption}
\label{ass:task}
(i) The object has one coordinate $\alpha\in[\alpha_{\min},\alpha_{\max}]$
and is always stably placed, so $\CP=\CS$.
(ii) Holding the object couples $\alpha$ to one robot joint $w$, the
\emph{wrist}: during a transfer, $\dot\alpha=-\dot w$. The sign follows from
the orientation of the axes.
(iii) Neither the grasp constraints nor collisions depend on $\alpha$:
$\CG=G\times[\alpha_{\min},\alpha_{\max}]$ and
$\CSfree=\Xfree\times[\alpha_{\min},\alpha_{\max}]$ for some
$G,\Xfree\subseteq\mathcal{X}$.
\end{assumption}
Point (iii) holds for rotationally symmetric objects, such as the cap in our
experiments. By (i), every grasping configuration is also a placement, so the
transition sets reduce to $G\cap\Xfree$ (times the object range). Transfer
motions also happen inside them.

The wrist has a limited range $[w_{\min},w_{\max}]$, so turning the object by
more than this range requires several grasps. Without loss of generality, let
$\agoal\le\ainit$. Since $\alpha$ affects neither collisions nor grasps, a plan
never gains from turning the object backwards, and the task works like a
ratchet. A \emph{stroke} grasps the object and moves the wrist forward, so
$\alpha$ decreases by the wrist travel. A \emph{return} releases the object
and rotates the wrist back with $\alpha$ fixed, ready for the next stroke. A
\emph{release} lets go of the object before the robot leaves the transition
set. Every other motion is transit in $\Xfree$ with $\alpha$ fixed.

\subsection{Problem statement}
\label{sec:problem:statement}
Given $\qinit=(\xinit,\ainit)$ and $\qgoal=(\xgoal,\agoal)$, find a continuous
path from $\qinit$ to $\qgoal$ in $\CSfree$, made of transit segments and
strokes, that uses the minimum number of grasps. Ties are broken by the number
of changes of convex set along the path. Path length is not optimised.

In terms of reachability, let $\mathcal{R}_k\subseteq\CSfree$ be the set of
states reachable from $\qinit$ with at most $k$ grasps. The minimum number of
grasps is the smallest $k$ such that $\qgoal\in\mathcal{R}_k$, and the goal
is unreachable if there is no such $k$. Neither $\mathcal{R}_k$ nor $\CSfree$
can be computed exactly. Our planner computes $\mathcal{R}_k$ exactly for an
inner approximation of $\Xfree$ and $G$ by convex polytopes, and its
guarantees hold relative to that approximation.
